\subsection{Identity and Inverse} 
Here are some properties that only hold if we have ``big enough'' number sets.
\begin{displaybox}[Identity]
There is a number \makeblanksmall the \term{additive identity} such that, for any real number $a$,
\[
\makeblanksmall+a=a+\makeblanksmall=a
\]
There is also a number \makeblanksmall the \term{multiplicative identity} such that, for any real number $a$,
\[
\makeblanksmall\cdot a=a\cdot\makeblanksmall=a
\]
\end{displaybox}
So far, we've only talked about addition and multiplication, not subtraction and division. Why?
\begin{displaybox}[Inverse]
For any real number $a$, there is a number \makeblanksmall called its \term{additive inverse} such that
\[
a+\makeblanksmall = \makeblanksmall+a = 0
\]
For any real number $a$ \emph{except for zero}, there is a number \makeblanksmall called its \term{multiplicative inverse} such that
\[
a\cdot\makeblanksmall = \makeblanksmall\cdot a=1
\]
\end{displaybox}
\begin{Note}
We wrote the multiplicative inverse of a number $a$ as $\frac{1}{a}$, but in general, to find the multiplicative inverse, we take the \makeblank[1.5in] of the number.
\begin{itemize}
\item The multiplicative inverse of $4$ is \makeblanksmall
\item The multiplicative inverse of $\frac{2}{3}$ is \makeblanksmall
\end{itemize}
\end{Note}
\begin{definition}[Subtraction and Division]
We define a new operation, \term{subtraction}, based on addition, by
\[
a-b=\makeblanksmall + \makeblanksmall
\]
We define a new operation, \term{division}, based on multiplication, by
\[
a\div b=\makeblanksmall\cdot\makeblanksmall
\]
\end{definition}
